Let $\mathfrak c$ denote the initial ordinal of cardinality $2^{\aleph_0}$ (the cardinality of $\mathbb R$). Let $\beta$ be a countable ordinal and let $n$ be an integer with $2\le n<\omega$.

For an ordinal $\kappa$ and integer $r\ge 1$, write $[\kappa]^r$ for the set of $r$-element subsets of $\kappa$.
For $\kappa,\beta$ ordinals and $n\in\mathbb N$, the partition relation
\[\kappa\to(\beta,n)_2^3\]
means:

\emph{for every} coloring $c:[\kappa]^3\to\{0,1\}$, \emph{either}
\begin{itemize}
\item there exists $H\subseteq\kappa$ of order type $\beta$ such that $c$ is constantly $0$ on $[H]^3$, \emph{or}
\item there exists $K\subseteq\kappa$ with $|K|=n$ such that $c$ is constantly $1$ on $[K]^3$.
\end{itemize}

Question: is it true that $\mathfrak c\to(\beta,n)_2^3$ for every countable ordinal $\beta$ and every integer $2\le n<\omega$?
